\documentclass[a4paper,12pt]{article}

% Pacchetti fondamentali
\usepackage[utf8]{inputenc}   % Codifica caratteri
\usepackage[italian]{babel}    % Lingua italiana
\usepackage{amsmath}           % Formule matematiche avanzate
\usepackage{graphicx}          % Inserimento immagini
\usepackage{geometry}          % Gestione margini
 \geometry{margin=2.5cm}
\usepackage{booktabs}          % Tabelle professionali
\usepackage{physics}
\usepackage{float}
\usepackage{caption}
\usepackage{textcomp}
\usepackage{amssymb}
\usepackage{parskip}
\usepackage{float}

\title{Misura del calore specifico a bassa temperatura}
\author{Eddi Bressan, Davide Conforte}
\date{\today}

\begin{document}

\maketitle
\tableofcontents
\newpage


\section{Obiettivo dell'esperienza}

L'obiettivo dell'esperienza è la determinazione sperimentale del calore latente di evaporazione
dell'azoto liquido e la successiva misura del calore specifico medio a pressione costante di due
campioni solidi: grafite e piombo.


\section{Cenni teorici}
\subsection{Calore latente di vaporizzazione dell'azoto}

L'azoto liquido assorbe calore $\Delta Q_0$ dall'ambiente e quindi evapora con una velocità
$\frac{\Delta m }{\Delta t}$ che, se le condizioni di contatto termico con l'ambiente rimangono invariate,
si può dire costante:
\begin{equation*}
    W_0 = \frac{\Delta Q_0}{\Delta t} = \frac{\lambda_f \Delta m}{\Delta t}
\end{equation*}
dove $\lambda_f$ è il calore latente di vaporizzazione dell'azoto.
Da cui:
\begin{equation*}
    \Delta m = \frac{W_0}{\lambda_f} \Delta t
\end{equation*}
Quando si fa passare corrente nel circuito con la resistenza a contatto con l'azoto liquido, la
quantità di calore assorbita aumenta di una frazione corrispondente alla potenza dissipata dalla
resistenza $W_R$:
\begin{equation*}
    W_0 + W_R = \frac{\Delta Q_0 + \Delta Q_R}{\Delta t}
\end{equation*}
Da cui:
\begin{equation*}
    \Delta m = \frac{W_0 + W_R}{\lambda_f} \Delta t
\end{equation*}
Le pendenze delle rette con il circuito aperto (senza corrente SC) e con il circuito chiuso
(con corrente CC) corrispondono alla variazione di massa nell'unità di tempo:
\begin{align*}
    \left. \frac{\Delta m}{\Delta t} \right|_{SC} &= \frac{W_0}{\lambda_f} \\
    \left. \frac{\Delta m}{\Delta t} \right|_{CC} &= \frac{W_0}{\lambda_f} + \frac{W_R}{\lambda_f}
\end{align*}

Sottraendo alla pendenza misurata con la corrente quella misurata senza, posso isolare il solo
contributo della resistenza:
\begin{equation*}
    \left. \frac{\Delta m}{\Delta t} \right|_{CC} - \left. \frac{\Delta m}{\Delta t} \right|_{SC} = \frac{W_R}{\lambda_f} = \frac{V I}{\lambda_f}
\end{equation*}
dove V e I sono rispettivamente la differenza di tensione e l'intensità di corrente misurate con
l'amperometro e il voltmetro (in serie e in parallelo con la resistenza).

In questo modo si riesce a ottenere una misura del calore latente di vaporizzazione dell'azoto:
\begin{equation*}
    \lambda_f = \frac{V I}{\left. \frac{\Delta m}{\Delta t} \right|_{CC} - \left. \frac{\Delta m}{\Delta t} \right|_{SC}}
\end{equation*}
con errore:
\begin{align*}
    \Delta \lambda_f &= \Delta V \left| \frac{I}{\left. \frac{\Delta m}{\Delta t} \right|_{CC} - \left. \frac{\Delta m}{\Delta t} \right|_{SC}} \right| + \Delta I \left| \frac{V}{\left. \frac{\Delta m}{\Delta t} \right|_{CC} - \left. \frac{\Delta m}{\Delta t} \right|_{SC}} \right| + \\
    &\quad + \left[ \Delta\left(\left. \frac{\Delta m}{\Delta t} \right|_{CC}\right) + \Delta \left( \left. \frac{\Delta m}{\Delta t} \right|_{SC} \right) \right] \left| \frac{VI}{\left( \left. \frac{\Delta m}{\Delta t} \right|_{CC} - \left. \frac{\Delta m}{\Delta t} \right|_{SC} \right)^2} \right|
\end{align*}

In alternativa possiamo usare il metodo grafico.
Indicando con $t_{start}$ e $t_{stop}$ gli istanti di accensione e di spegnimento, si tracciano le
rette per la fase 1 (prima dell'accensione) e per la fase 3 (dopo l'accensione) e le si estrapola
a $t_m = \frac{t_{start} + t_{stop}}{2}$.

Poi si tracciano le rette a pendenza massima e minima, estrapolandole allo stesso $t_m$ e
si misurano i valori di $\Delta \left( \left.\Delta m \right|_{fase 1} \right)$ e $\Delta \left( \left.\Delta m \right|_{fase 3} \right)$.

Si misurano quindi la distanza tra le due rette al tempo $t_m$ di valore $\Delta m$ e l'errore $\Delta \left( \Delta m \right) = \Delta \left( \left.\Delta m \right|_{fase 1} \right) + \Delta \left( \left.\Delta m \right|_{fase 3} \right)$.

Si può così misurare il calore latente di vaporizzazione:
\begin{equation*}
    \lambda_f = \frac{V I}{\Delta m} (t_{stop} - t_{start}) = \frac{V I \Delta t}{\Delta m}
\end{equation*}
con errore:
\begin{equation*}
    \Delta \lambda_f = \lambda_f \left( \frac{\Delta V}{V} + \frac{\Delta I}{I} + \frac{\Delta (\Delta m)}{(\Delta m)} + \frac{\Delta (\Delta t)}{(\Delta t)} \right)
\end{equation*}
Per questo metodo, si possono ricavare le equazioni delle rette con il metodo dei minimi quadrati
lineari per trovare il migliore fit lineare.


\subsection{Calore specifico medio}

Il calore trasferito da n moli di un materiale all'azoto liquido ad una certa temperatura è:
\begin{equation*}
    dQ = nc(T)dT
\end{equation*}
con c calore specifico del materiale alla temperatura T.

Nell'intervallo tra la temperatura ambiente $T_a$ e la temperatura di vaporizzazione dell'azoto
liquido $T_f$ si ha:
\begin{equation*}
    Q = n \int_{T_a}^{T_f} c(T)dT \approx n \sum \frac{c(T_k) + c(T_{k+1})}{2} \abs{T_{k+1} - T_k}
\end{equation*}
Si definisce il calore specifico medio nell'intervallo di temperatura come:
\begin{equation*}
    \bar{c_x} = \frac{1}{n} \frac{Q}{\abs{T_a - T_f}} = \frac{1}{\abs{T_a - T_f}} \sum \frac{c(T_k) + c(T_{k+1})}{2} \abs{T_{k+1} - T_k}
\end{equation*}
Formula che può essere utilizzata per il calcolo approssimato degli integrali (regola del trapezio)
per calcolare la capacità termica molare media per le sostanze di nostro interesse.

Pertanto, usando il metodo grafico e estrapolando le rette per la fase 1 (prima dell'immersione
della sostanza) e per la fase 3 (dopo la stabilizzazione della perdita di massa nel tempo) come
descritto sopra, si può calcolare il calore specifico medio (J/(Kg K)) come:
\begin{equation*}
        \bar{c_x} = \frac{\Delta m \lambda_f}{m_x (T_f - T_a)}
\end{equation*}
con $m_x$ la massa della sostanza. L'errore di $c_x$ è dato da:
\begin{equation*}
    \Delta c_x = \left( \frac{\Delta (\Delta m)}{(\Delta m)} + \frac{\Delta \lambda_f}{\lambda_f} + \frac{\Delta m_x}{m_x} + \frac{2 \Delta T}{(T_a - T_f)} \right) c_x
\end{equation*}
Di seguito si riporta una tabella di calori specifici molari a diverse temperature, su cui possiamo basarci per 
fare predizioni utili allo svolgimento della esperienza:
\begin{table}[H]
\centering
\begin{tabular}{ccc}
\toprule
\textbf{T} (K) & \textbf{C} ($\text{cal}\cdot\text{mol}^{-1}\cdot\text{K}^{-1}$) & \textbf{Pb} ($\text{cal}\cdot\text{mol}^{-1}\cdot\text{K}^{-1}$) \\
\midrule
70  & 0.221 & 5.54 \\
80  & 0.278 & 5.64 \\
90  & 0.339 & 5.74 \\
100 & 0.402 & 5.84 \\
120 & 0.539 & 5.94 \\
140 & 0.689 & 5.99 \\
160 & 0.849 & 6.09 \\
180 & 1.019 & 6.14 \\
200 & 1.188 & 6.19 \\
220 & 1.36  & 6.24 \\
240 & 1.54  & 6.29 \\
260 & 1.71  & 6.34 \\
280 & 1.88  & 6.38 \\
300 & 2.05  & 6.44 \\
\bottomrule
\end{tabular}
\end{table}




\section{Apparato sperimentale}

Per la misura del calore latente di vaporizzazione è stato realizzato un circuito composto da:
\begin{itemize}
    \item Resistenza elettrica da $10.3 \pm 0.1 \Omega$
    \item Generatore di corrente continua
    \item Voltmetro con errore di sensibilità ....DA INSERIRE....
    \item Amperometro con errore di sensibilità ....DA INSERIRE....
\end{itemize}

Sono stati utilizzati anche i seguenti strumenti e materiali:
\begin{itemize}
    \item Vaso Dewar
    \item Azoto liquido ($LN_2$)
    \item Campione di piombo ($Pb$)
    \item Campione di grafite ($C$)
    \item Termometro a termocoppia con errore di sensibilità ....DA INSERIRE....
    \item Bilancia elettronica con errore di sensibilità ....DA INSERIRE....
    \item Cronometro digitale con errore di sensibilità ....DA INSERIRE....
\end{itemize}




\section{Procedura sperimentale}

Come procedimento preliminare all'esperienza, è stata calcolata la potenza dissipata dal circuito necessaria
a indurre l'evaporazione di una massa di circa 10g di azoto liquido in un intervalllo di 5 minuti.
Il generatore di corrente è stato conseguentemente impostato sui valori di tensione e intensità di
corrente atti a garantire tale potenza.

Successivamente, misurata la temperatura ambiente, si è proceduto alla determinazione del calore 
specifico medio dei campioni di piombo e grafite, utilizzando i dati riportati nella tabella in
precedenza.

Infine, attraverso un bilancio termido basato sulla temperatura di equilibrio raggiunta, è stata
stimata la quantità di azoto evaporato in seguito al contatto con i materiali analizzati.

\subsection{Calore latente di vaporizzazione dell'azoto}
Si inserisce l'azoto liquido nel vaso Dewar insieme alla resistenza elettrica e si posiziona il tutto
sopra la bilancia elettronica.

Si misura così la massa del sistema ogni 20 secondi per 5 minuti. Terminati i 5 minuti si accende il circuito
e si continua a misurare la massa ogni 10 secondi per 5 minuti. Infine si spegne il circuito e si continua
a misurare la massa ogni 20 secondi per 5 minuti.


\subsection{Calore specifico medio}

Il seguente procedimento va eseguito per il campione di piombo e quello di grafite.

Si misura la massa del campione, poi si posiziona il vaso Dewar con l'azoto liquido sulla bilancia
elettronica insieme al campione che però resta posizionato al di fuori del vaso Dewar, sopra il piatto della bilancia.

Si misura così la massa del sistema ogni 20 secondi per 5 minuti. Si introduce poi il campione nell'azoto
liquido e si continua a misurare la massa del sistema ogni 10 secondi finché il campione e l'azoto liquido
non raggiungono l'equilibrio termico tra loro e si registra una perdita di massa dovuta solamente allo scambio
termico con l'ambiente.

\section{Presa dati}
In seguito sono riportate tutte le misure raccolte durante l'esperienza.

Per avere una perdita pari a 10g di azoto in 5 minuti, è stato calcolata una potenza necessaria di:
\begin{equation*}
    W = VI = 6.67W
\end{equation*}
Pertanto è stato misurato:

\begin{table}[h]
    \centering
    \begin{tabular}{|c|c|} \hline
        V (V) & 8.29 \\ \hline
        I (A) & 0.80 \\ \hline
    \end{tabular}
\end{table}

Le masse misurate dei campioni sono:

\begin{table}[h]
    \centering
    \begin{tabular}{|c|c|} \hline
        Piombo (g) & 38.48 \\ \hline
        Grafite (g) & 17.17 \\ \hline
    \end{tabular}
\end{table}

Ricaviamo così il numero di moli:
\begin{table}[h]
    \centering
    \begin{tabular}{|c|c|c|c|} 
        \hline
        \multicolumn{2}{|c|}{\textbf{Piombo}} & \multicolumn{2}{c|}{\textbf{Grafite}} \\ \hline
        M (g/mol) & 207.2 & M (g/mol) & 12.01 \\ \hline
        n (mol) & 1.430 & n (mol) & 0.186 \\ \hline
    \end{tabular}
\end{table}

Avendo misurato $T_a = 294.6K$, sapendo che $T_f = 77.4K$ e avendo supposto lineare la variazione di calore
specifico negli intervalli indicati nella tabella, abbiamo calcolato il calore specifico medio e la perdita
di massa: 


DAVIDE AIUTO!!!!!!!!!!!!!!
\begin{table}[H]
    \centering
    % Prima tabella
    \begin{minipage}{0.45\textwidth}
        \centering
        \begin{tabular}{|c|c|}
            \hline
            \multicolumn{2}{|c|}{\textbf{Piombo}} \\ \hline
            $C_x$ (cal/(mol K)) & Dato 2 \\ 
            $C_x$ (J/(Kg K)) & Dato 4 \\ \hline
            $\Delta m$ (g) & Dato 6 \\ \hline
        \end{tabular}
        \caption{Prima tabella}
    \end{minipage}
    \hfill
    % Seconda tabella
    \begin{minipage}{0.45\textwidth}
        \centering
        \begin{tabular}{|c|c|}
            \hline
            \multicolumn{2}{|c|}{\textbf{Grafite}} \\ \hline
            $C_x$ (cal/(mol K)) & Dato 2 \\ 
            $C_x$ (J/(Kg K)) & Dato 4 \\ \hline
            $\Delta m$ (g) & Dato 6 \\ \hline
        \end{tabular}
        \caption{Seconda tabella}
    \end{minipage}
\end{table}
DAVIDE AIUTO!!!!!!!!!!!!!
Abbiamo quindi una perdita di massa totale XYZ, escludendo quella continuamente persa a causa dello scambio
termico con l'ambiente.


\subsection{Calore latente di vaporizzazione dell'azoto}

\begin{table}[H]
    \centering
    % Prima tabella
    \begin{minipage}[t]{0.30\textwidth}
        \vspace{0pt}
        \centering
        \begin{tabular}{|c|c|}
            \hline
            \multicolumn{2}{|c|}{\textbf{Senza corrente}} \\ \hline
            t (s) & m (g) \\ \hline
            0.00 & 1606.58 \\ \hline
            20.35 & 1606.20 \\ \hline
            40.41 & 1605.65 \\ \hline
            60.11 & 1605.25 \\ \hline
            80.62 & 1604.75 \\ \hline
            100.93 & 1604.30 \\ \hline
            122.01 & 1603.87 \\ \hline
            142.22 & 1603.38 \\ \hline
            162.67 & 1602.96 \\ \hline
            184.60 & 1602.41 \\ \hline
            204.96 & 1602.00 \\ \hline
            225.69 & 1601.55 \\ \hline
            245.66 & 1601.06 \\ \hline
            265.77 & 1600.65 \\ \hline
            289.19 & 1600.16 \\ \hline
            306.70 & 1599.70 \\ \hline
        \end{tabular}
    \end{minipage}
    \hfill
    % Seconda tabella
    \begin{minipage}[t]{0.30\textwidth}
        \vspace{0pt}
        \centering
        \begin{tabular}{|c|c|}
            \hline
            \multicolumn{2}{|c|}{\textbf{Con corrente}} \\ \hline
            t (s) & m (g) \\ \hline
            365.48 & 1598.70 \\ \hline
            375.86 & 1598.30 \\ \hline
            386.20 & 1597.50 \\ \hline
            396.84 & 1597.00 \\ \hline
            406.82 & 1596.50 \\ \hline
            417.29 & 1596.00 \\ \hline
            427.53 & 1595.40 \\ \hline
            437.65 & 1594.25 \\ \hline
            456.31 & 1593.77 \\ \hline
            466.67 & 1593.12 \\ \hline
            477.13 & 1592.70 \\ \hline
            487.43 & 1592.09 \\ \hline
            497.75 & 1591.50 \\ \hline
            507.99 & 1590.99 \\ \hline
            518.56 & 1590.40 \\ \hline
            528.98 & 1589.78 \\ \hline
            539.90 & 1589.30 \\ \hline
            550.28 & 1588.80 \\ \hline
            560.86 & 1588.27 \\ \hline
            571.96 & 1587.62 \\ \hline
            582.27 & 1587.08 \\ \hline
            592.65 & 1586.59 \\ \hline
            603.03 & 1585.98 \\ \hline
        \end{tabular}
    \end{minipage}
    \hfill
    % Terza tabella
    \begin{minipage}[t]{0.30\textwidth}
        \vspace{0pt}
        \centering
        \begin{tabular}{|c|c|}
            \hline
            \multicolumn{2}{|c|}{\textbf{Senza corrente}} \\ \hline
            t (s) & m (g) \\ \hline
            623.40 & 1585.16 \\ \hline
            643.85 & 1584.90 \\ \hline
            664.21 & 1584.70 \\ \hline
            684.70 & 1584.40 \\ \hline
            705.06 & 1584.94 \\ \hline
            725.79 & 1584.45 \\ \hline
            746.03 & 1584.02 \\ \hline
            766.48 & 1583.63 \\ \hline
            787.24 & 1583.25 \\ \hline
            807.74 & 1582.80 \\ \hline
            828.15 & 1582.40 \\ \hline
            848.36 & 1582.96 \\ \hline
            868.78 & 1581.55 \\ \hline
            889.21 & 1581.13 \\ \hline
        \end{tabular}
    \end{minipage}
\end{table}


\subsection{Calore specifico medio del piombo}






\section{Analisi dati}

\section{Conclusioni}

\end{document}